

\paragraph{Hyperparameters.}

We tuned the BERT model with a grid search over the following hyperparameters: learning rate (1e-5, 2e-5, 3e-5), sequence length (96, 128, 160), and batch size [16, 32], this resulted in 18 configurations. We selected the best configuration based on macro-F1 on the development set.

The LSTM model was manually tuned using five hand-designed experiments each isolating one hyperparameter (early stopping, regularization, sequence length, data scale). We selected the best hyperparameters based on macro-F1 on the development set.

We report accuracy and macro-F1 as our main evaluation metrics, with macro-F1 used as the primary measure for model selection because it gives equal weight to all four classes. In addition, we generated classification reports, confusion matrices, learning curves, and files with misclassified test examples for later error analysis. The test set was only used for final evaluation after model selection.

\paragraph{Shared training settings.}
Both models were trained with the Adam optimiser. We used gradient clipping with a maximum norm of 1.0 and early stopping with a patience of 2 epochs based on validation loss. The batch size for BERT was set to 16, while for the LSTM it was set to 32 to boost the training times.

\paragraph{Robustness/Slicing Evaluation: Length buckets.}
We evaluated the models on different length buckets to see how performance varies with input length. We created four buckets based on the distribution of word counts in the training set: shortest (0-34 words), short (35-39 words), long (40-45 words), and longest (46+ words). We then evaluated the models separately on the test examples that fall into each bucket. This analysis allows us to understand if the models perform better on shorter or longer texts, which can reveal insights about their ability to capture context and handle varying input lengths. For instance, if a model performs significantly worse on long texts, it may indicate that it struggles to capture long-range dependencies or that it is overfitting to shorter patterns in the training data.

\paragraph{Robustness/Slicing Evaluation: Keyword Masking Probe.}
We implemented a keyword masking probe to test how much the models rely on specific keywords for classification. We identified the top 6 most informative keywords per class (24 in total) based on the training data using TF-IDF scores. Then, we created a list of control words, from the whole training vocabulary, that have the same frequency distribution as the keywords but are not informative for classification (i.e., have the same TF score but low IDF). Then, we created a modified test set where these keywords were masked out (replaced with a special [MASK] token), we did the same for the control words, and we evaluated the models on both masked test sets to see how much their performance drops compared to the original test set. A significant drop in performance on the keyword-masked test set would suggest that the model relies heavily on these keywords for classification, while a smaller drop would indicate that the model is using more contextual information and is less sensitive to specific keywords. The control word masking serves as a baseline to ensure that any observed performance drop is indeed due to the masking of informative keywords rather than just any random words. This analysis helps us understand the robustness of the models to changes in key lexical cues and whether they are learning more generalizable patterns or just memorizing specific terms.


% Then, we created a modified test set where these keywords were masked out (replaced with a special [MASK] token). We evaluated the models on this masked test set to see how much their performance drops compared to the original test set. A significant drop in performance would suggest that the model relies heavily on these keywords for classification, while a smaller drop would indicate that the model is using more contextual information and is less sensitive to specific keywords. This analysis helps us understand the robustness of the models to changes in key lexical cues and whether they are learning more generalizable patterns or just memorizing specific terms.

